\documentclass{article}
\usepackage{fontspec}
\setmainfont{Libertinus Serif}
\usepackage[a4paper, margin=1in]{geometry}
\usepackage{amsmath, amsthm, amssymb, mathtools, enumitem, hyperref}
\usepackage[portuguese]{babel}

\setlist[description]{leftmargin=1cm,labelindent=1cm}

\theoremstyle{plain}
\newtheorem{thm}{Teorema}[section]
\newtheorem{prop}[thm]{Proposição}
\newtheorem{lem}[thm]{Lema}
\newtheorem{cor}[thm]{Corolário}
\newtheorem{res}[thm]{Resultado}

\theoremstyle{definition}
\newtheorem{defn}{Definição}[section]
\newtheorem{exmp}[defn]{Exemplo}

\theoremstyle{remark}
\newtheorem*{rem}{Nota}

\begin{document}
\begin{center}
    \textbf{\large UNIVERSIDADE FEDERAL DO RIO GRANDE DO SUL} \\
    \textbf{\large INSTITUTO DE MATEMÁTICA E ESTATÍSTICA} \\
    \vspace{3cm}
    \textbf{\Large MÉTODOS APLICADOS DE MATEMÁTICA II} \\
    \vspace{1cm}
    \textbf{\large LISTA DE DEFINIÇÕES E RESULTADOS} \\
    \vspace{3cm}

    \textbf{LEONARDO SUAREZ}

    \vfill
    \textbf{\large Porto Alegre, RS} \\
    \textbf{\large \today}
\end{center}

\newpage
\tableofcontents
\newpage

\section{Soluções por Séries de Potências}

\subsection{Revisão de Séries de Potências}

\begin{rem} (fatos importantes sobre séries de potências).
    \begin{itemize}
        \item \textbf{Definição.} Uma série de potências centrada em \(x-a\) é uma série infinita da forma \begin{equation}\sum_{n=0}^{\infty} c_n (x-a)^n = c_0 + c_1(x-a) + c_2(x-a)^2 + \cdots\end{equation}
        \item \textbf{Convergência.} Uma série de potências é convergente no ponto \(x=x_0\) se a sua sequência de somas parciais \(\{S_N(x)\}\) converge para um número real para \(N\to\infty\), ou seja, \(\lim_{N\to\infty} S_N(x_0) = \lim_{N\to\infty} \sum_{n=0}^{N} c_n (x-a)^n\) existe. Se o limite não existe, a série diverge em \(x_0\).
        \item \textbf{Intervalo de convergência.} O conjunto de todos os pontos \(x\) para os quais a série de potências converge é chamado de intervalo de convergência. O centro do intervalo de convergência é o centro \(a\) da série de potências.
        \item \textbf{Raio de convergência.} O raio \(R\) do intervalo de convergência de uma série de potências é chamado de raio de convergência. Se a série converge para todo \(x\), então \(R=\infty\). Se \(R>\), então a série converge para \(|x-a|<R\) e diverge para \(|x-a|>R\). Se \(R=0\), então a série converge somente em \(x=a\). Uma série de potências pode ou não convergir nos pontos finais do intervalo de convergência.
        \item \textbf{Convergência absoluta.} Dentro de seu intervalo de convergência, uma série de potências converge absolutamente, ou seja, a série \(\sum_{n=0}^{\infty} |c_n (x-a)^n|\) também converge para todo \(x \in (a-R, a+R)\).
        \item \textbf{Teste da razão.} A convergência de uma série de potências pode ser determinada pelo teste da razão. Suponha que \(c_n\neq0\) para todo \(n\) em \(\sum_{n=0}^{\infty} c_n (x-a)^n\), e que 
        \begin{equation}
            \lim_{n\to\infty} \left| \frac{c_{n+1}(x-a)^{n+1}}{c_n (x-a)^n} \right| = |x-a| \lim_{n\to\infty} \left| \frac{c_{n+1}}{c_n} \right| = L.
        \end{equation}
        \noindent Se \(L<1\), a série converge absolutamente; se \(L>1\), a série diverge; se \(L=1\), o teste é inconclusivo. O teste da razão é sempre inconclusivo nos pontos \(x=a\pm R\).
        \item \textbf{Uma série de potências define uma função.} Uma série de potências \(\sum_{n=0}^{\infty} c_n (x-a)^n\) define uma função \(f\) por \( f(x) = \sum_{n=0}^{\infty} c_n (x-a)^n\) para todo \(x\) no intervalo de convergência. Se o raio de convergência é \(R>0\) ou \(R = \infty\), então \(f\) é contínua, diferenciável e integrável nos intervalos \((a-R, a+R)\) ou \((-\infty, \infty)\), respectivamente. Além disso, \(f\) pode ser diferenciada e integrada termo a termo. Convergência em um ponto final do intervalo pode ser perdida após a diferenciação ou integração. Se 
        \[ f(x) = \sum_{n=0}^{\infty} c_n x^n = c_0 + c_1 x + c_2 x^2 + \cdots \]
        \noindent é uma série de potências em \(x\), então suas primeiras duas derivadas são \( y' = \sum_{n=1}^{\infty} n c_n x^{n-1}\) e \( y'' = \sum_{n=2}^{\infty} n(n-1)c_n x^{n-2}\).
        \item \textbf{Identidade.} Se \(\sum_{n=0}^{\infty} c_n (x-a)^n = 0\) para todo \(x\) em um intervalo, então \(c_n = 0\) para todo \(n\).
        \item \textbf{Analítica em um ponto.} Uma função \(f\) é dita analítica em um ponto \(x=a\) se existe uma série de potências centrada em \(x-a\) que converge para \(f(x)\) em algum intervalo \((a-R, a+R)\), onde \(R>0\).
        \item \textbf{Aritmética de séries de potências.} Séries de potências podem ser combinadas pelas operações de soma, multiplicação e divisão. Os procedimentos são semelhantes aos usados para polinômios.
    \end{itemize}
\end{rem}

\subsection{Soluções sobre Pontos Ordinários}

\begin{defn}
    Um ponto \(x=x_0\) é chamado de \emph{ponto ordinário} da equação diferencial
    \begin{equation}\label{eq:edo1}
        a_2(x)y'' + a_1(x)y' + a_0(x)y = 0,
    \end{equation}
    \noindent se ambos os coeficientes \(P(x)\) e \(Q(x)\) na forma padrão
    \begin{equation}\label{eq:edo2}
        y'' + P(x)y' + Q(x)y = 0,
    \end{equation}
\noindent são analíticos em \(x_0\). Caso contrário, \(x_0\) é chamado de \emph{ponto singular}.
\end{defn}

\begin{thm} (existência de soluções por séries de potências).
    Se \(x_0\) é ponto ordinário da equação \eqref{eq:edo1}, então podemos encontrar duas soluções linearmente independentes na forma de séries de potências centradas em \(x_0\), ou seja,
    \[
    y = \sum_{n=0}^{\infty} c_n (x - x_0)^n.
    \]
    Uma solução de série de potência converge pelo menos no intervalo \(|x - x_0| < R\), onde \(R\) é a distância do ponto \(x_0\) ao ponto singular mais próximo, ou também a \emph{cota inferior} do raio de convergência.
\end{thm}

\subsection{Soluções sobre Pontos Singulares}

\begin{defn}
    Uma ponto singular \(x=x_0\) é chamado de \emph{ponto singular regular} da equação \eqref{eq:edo1} se as funções \(p(x)=(x-x_0)P(x)\) e \(q(x)=(x-x_0)^2Q(x)\) são analíticas em \(x_0\). Caso contrário, \(x_0\) é chamado de \emph{ponto singular irregular}.
\end{defn}

\begin{thm} (teorema de Frobenius).
    Se \(x_0\) é ponto singular regular da equação \eqref{eq:edo1}, então existe pelo menos uma solução da forma
    \begin{equation}
        y = (x - x_0)^r \sum_{n=0}^{\infty} c_n (x - x_0)^n = \sum_{n=0}^{\infty} c_n (x - x_0)^{n+r},
    \end{equation}
    \noindent onde \(r\) é a raiz de uma equação indicial associada. A série converge pelo menos em algum intervalo \(0<x-x_0<R\).
\end{thm}

\begin{rem}
    Suponha que \(x=0\) é ponto singular regular da equação \eqref{eq:edo1} e que as raízes \(r_1\) e \(r_2\) da equação indicial associada satisfaçam \(r_1 \geq r_2\). Então:
    \begin{itemize}
        \item [\textbf{Caso I.}] Se \(r_1 - r_2\) não é inteiro positivo, então existem duas soluções linearmente independentes da forma
        \begin{equation}
            y_1 = x^{r_1} \sum_{n=0}^{\infty} c_n x^n, \quad c_0 \neq 0, \qquad y_2 = x^{r_2} \sum_{n=0}^{\infty} d_n x^n, \quad d_0 \neq 0.
        \end{equation}
        \item [\textbf{Caso II.}] Se \(r_1 - r_2 = N\), onde \(N\) é um inteiro positivo, então existem soluções da forma
        \begin{equation}
            y_1 = x^{r_1} \sum_{n=0}^{\infty} c_n x^n, \quad c_0 \neq 0, \qquad y_2 = C y_1 \ln x + x^{r_2} \sum_{n=0}^{\infty} d_n x^n, \quad d_0 \neq 0,
        \end{equation}
        \noindent onde \(C\) é uma constante que pode ser zero.
        \item [\textbf{Caso III.}] Se \(r_1 = r_2\), então as soluções são da forma
        \begin{equation}
            y_1 = x^{r_1} \sum_{n=0}^{\infty} c_n x^n, \quad c_0 \neq 0, \qquad y_2 = y_1 \ln x + x^{r_2} \sum_{n=1}^{\infty} d_n x^n.
        \end{equation}
    \end{itemize}
\end{rem}

\subsection{Funções Especiais}

\begin{defn}
    Assuma que \(\nu\geq0\). Então,
    \begin{equation}\label{eq:bessel}
        x^2y'' + xy' + (x^2 - \nu^2)y = 0
    \end{equation}
    \noindent é chamada de \emph{equação de Bessel de ordem \(\nu\)}.
\end{defn}

\begin{res}
    A solução geral da equação de Bessel de ordem \(\nu\) é dada por
    \begin{equation}
        y = c_1 J_{\nu}(x) + c_2 J_{-\nu}(x), \quad \nu \notin \mathbb{Z},
    \end{equation}
    \noindent onde \(J_{\nu}(x)\) é a função de Bessel de primeira espécie e ordem \(\nu\), definida por
    \begin{equation}
        J_{\nu}(x) = \sum_{n=0}^{\infty} \frac{(-1)^n}{n! \Gamma(n+\nu+1)} \left(\frac{x}{2}\right)^{2n+\nu},
    \end{equation}
    \noindent que converge no intervalo \((0,\infty)\). O termo \(\Gamma(n+\nu+1)\) é a função gama, também conhecida como fatorial generalizado, definida por \(\Gamma(x) = \int_0^{\infty} t^{x-1} e^{-t} \, dt\) para \(x>0\) e satisfaz a propriedade \(\Gamma(x+1) = x\Gamma(x)\).
\end{res}

\begin{res}
    Para todo valor de \(\nu\), a solução geral de \eqref{eq:bessel} pode ser escrita como
    \begin{equation}
        y = c_1 J_{\nu}(x) + c_2 Y_{\nu}(x),
    \end{equation}
    \noindent onde \(Y_{\nu}(x)\) é a função de Bessel de segunda espécie e ordem \(\nu\), definida por
    \begin{equation}
        Y_{\nu}(x) = \frac{J_{\nu}(x) \cos(\nu \pi) - J_{-\nu}(x)}{\sin(\nu \pi)},
    \end{equation}
    \noindent que também converge no intervalo \((0,\infty)\).
\end{res}

\begin{prop} (algumas propriedades das funções de Bessel).
    Para funções de Bessel de ordem \(m\), \(m=0,1,2,\ldots\), temos as seguintes propriedades:
    \begin{align*}
        &J_{-m}(x) = (-1)^m J_m(x), \quad & &J_m(-x) = (-1)^m J_m(x), \\
        &J_m(0) = \begin{cases}
            0, \quad m > 0 \\
            1, \quad m = 0
        \end{cases}, \quad & &\lim_{x\to 0^{+}} Y_m(x) = -\infty
    \end{align*}
\end{prop}

\begin{defn}
    Considerando \(n\) um inteiro não-negativo, temos 
    \begin{equation}\label{eq:legendre}
        (1-x^2)y'' - 2xy' + n(n+1)y = 0
    \end{equation}
    \noindent que é chamada de \emph{equação de Legendre de ordem \(n\)}.
\end{defn}

\begin{res}
    Os polinômios de Legendre \(P_n(x)\) são soluções particulares da equação de Legendre de ordem \(n\), dadas por
    \begin{align*}
        P_0(x) &= 1, & \quad P_1(x) &= x,\\
        P_2(x) &= \frac{1}{2}(3x^2 - 1), & \quad P_3(x) &= \frac{1}{2}(5x^3 - 3x),\\
        P_4(x) &= \frac{1}{8}(35x^4 - 30x^2 + 3), & \quad P_5(x) &= \frac{1}{8}(63x^5 - 70x^3 + 15x). 
    \end{align*}

    \noindent A fórmula de Rodrigues, que gera os polinômios de Legendre, é dada por
    \begin{equation}\label{eq:rodrigues}
        P_n(x) = \frac{1}{2^n n!} \frac{d^n}{dx^n} (x^2 - 1)^n, \quad n=0,1,2,\ldots
    \end{equation}
\end{res}

\begin{prop} (algumas propriedades dos polinômios de Legendre).
    \begin{align*}
        &P_n(-x) = (-1)^n P_n(x), &\quad & P_n(1) = 1, \quad & P_n(-1) = (-1)^n, \\
        &P_n(0) = 0, \quad n \text{ ímpar}, &\quad & P'_n(0) = 0, \quad n \text{ par} &
    \end{align*}
\end{prop}

\begin{res} (relação de recorrência dos polinômios de Legendre).
    \begin{equation}\label{eq:recorrencia-legendre}
        (n+1)P_{n+1}(x) - (2n+1)xP_n(x) + nP_{n-1}(x) = 0, \quad n=1,2,3,\ldots
    \end{equation}
\end{res}

\section{Séries de Fourier}

\subsection{Funções Ortogonais}

\begin{defn} 
    O \emph{produto interno} de duas funções \(f_1\) e \(f_2\) em um intervalo \([a,b]\) é o número
    \begin{equation}
        \langle f_1, f_2 \rangle = \int_a^b f_1(x)f_2(x) \, dx.
    \end{equation}
\end{defn}

\begin{defn}
    Duas funções \(f_1\) e \(f_2\) são ditas \emph{ortogonais} em \([a,b]\) se
    \begin{equation}
        \langle f_1, f_2 \rangle = 0.
    \end{equation}
\end{defn}

\begin{defn}
    Um conjunto de funções reais \(\{\phi_1(x), \phi_2(x), \ldots\}\) é \emph{ortogonal} em \([a,b]\) se
    \begin{equation}
        \langle \phi_m, \phi_n \rangle = \int_a^b \phi_m(x)\phi_n(x) \, dx = 0, \quad m \neq n.
    \end{equation}
\end{defn}

\begin{defn}
    A \emph{norma quadrada} de uma função \(\phi_n\) é \(\Vert \phi _n \Vert^2 = \langle \phi_n, \phi_n \rangle = \int_a^b \phi_n(x)^2 \, dx\), e a \emph{norma} é dada por \(\Vert \phi_n \Vert = \sqrt{\langle \phi_n, \phi_n \rangle} = \sqrt{\int_a^b \phi_n(x)^2 \, dx}\). Se \(\Vert \phi_n \Vert = 1\), então \(\{\phi_n\}\) é um conjunto \emph{ortonormal} no intervalo \([a,b]\).
\end{defn}

\subsection{Séries de Fourier}

\begin{defn}
    Um conjunto de funções reais \(\{\phi_0(x), \phi_1(x), \phi_2(x), \ldots\}\) é \emph{ortogonal com respeito a uma função peso} \(w(x)\) em \([a,b]\) se
    \begin{equation}
        \int_{a}^{b} w(x)\phi_m(x)\phi_n(x) \, dx = 0, \quad m \neq n.
    \end{equation}

    \noindent A série 
    \begin{equation}
    f(x)=\sum_{n=0}^{\infty} c_n \phi_n(x)\,, \quad c_n=\frac{\int_{a}^{b} f(x)w(x)\phi_n(x) \, dx}{\Vert \phi_n(x)\Vert^2}.
    \end{equation}
    \noindent é chamada de \emph{série de Fourier generalizada} de \(f\).
\end{defn}

\begin{rem}
    O conjunto \(\left\{1, \cos\frac{n\pi x}{p}, \sin\frac{m\pi x}{p} \right\}, n=1,2,3,\ldots, m=1,2,3,\ldots\) é ortogonal em \([-p,p]\) com respeito à função peso \(w(x)=1\).
\end{rem}

\begin{defn}
    A \emph{série de Fourier} de uma função \(f\) definida no intervalo \((-p,p)\) é dada por
    \begin{equation}\label{eq:serie-fourier}
        f(x) = \frac{1}{2}a_0 + \sum_{n=1}^{\infty} \left( a_n \cos\frac{n\pi x}{p} + b_n \sin\frac{n\pi x}{p} \right),
    \end{equation}
    onde os coeficientes \(a_n\) e \(b_n\) são dados por
    \begin{align}
        a_0 &= \frac{1}{p} \int_{-p}^{p} f(x) \, dx, \\
        a_n &= \frac{1}{p} \int_{-p}^{p} f(x) \cos\frac{n\pi x}{p} \, dx, \\
        b_n &= \frac{1}{p} \int_{-p}^{p} f(x) \sin\frac{n\pi x}{p} \, dx.
    \end{align}
\end{defn}

\begin{rem}
    A série de Fourier \eqref{eq:serie-fourier} pode ser escrita na forma complexa como
    \begin{equation}\label{eq:serie-fourier-complexa}
        f(x) = \sum_{n=-\infty}^{\infty} c_n e^{i n \pi x / p},
    \end{equation}
    \noindent onde
    \begin{equation}
        c_0 = \frac{a_0}{2}, \quad c_n = \frac{1}{2}(a_n - i b_n), \quad c_{-n} = \frac{1}{2}(a_n + i b_n), \quad n=1,2,3,\ldots
    \end{equation}
\end{rem}

\begin{defn}
    Uma função \(f\) é dita \emph{contínua por partes} em um intervalo \([a,b]\) se existe um número finito de pontos de descontinuidade \(x_1<x_2<\cdots<x_n\) em \([a,b]\) e \(f\) é contínua em cada subintervalo \((x_k, x_{k+1})\) para \(k=0,1,\ldots,n\). Consequentemente,
    \begin{itemize}
        \item Os limites laterais \(f(x+)\) e \(f(x-)\) existem para todo \(x \in (a,b)\).
        \item Os limites laterais \(f(a+)\) e \(f(b-)\) existem, mas não é necessário que \(f\) seja contínua ou mesmo definida em \(a\) ou em \(b\)
    \end{itemize}
\end{defn}

\begin{thm} (condições para convergência).
    Sejam \(f\) e \(f'\) contínuas por partes em \([-p,p]\). Então para todo \(x \in (-p,p)\), a série de Fourier de \(f\) converge para \(f(x)\) em um ponto de continuidade. Em um ponto de descontinuidade, a série de Fourier converge para a média
    \begin{equation*}
        \frac{f(x+) + f(x-)}{2},
    \end{equation*}
    \noindent onde \(f(x+)\) e \(f(x-)\) são os limites laterais de \(f\) em \(x\), respectivamente.
\end{thm}

\begin{thm} (propriedades de funções pares/ímpares).
    \begin{enumerate}
        \item[\textup{(a)}] O produto de duas funções pares é par.
        \item[\textup{(b)}] O produto de duas funções ímpares é par.
        \item[\textup{(c)}] O produto de uma função par por uma função ímpar é ímpar.
        \item[\textup{(d)}] A soma (diferença) de duas funções pares é par.
        \item[\textup{(e)}] A soma (diferença) de duas funções ímpares é ímpar.
        \item[\textup{(f)}] Se \(f\) é par, então \(\int_{-a}^{a} f(x) \, dx = 2\int_{0}^{a} f(x) \, dx\).
        \item[\textup{(g)}] Se \(f\) é ímpar, então \(\int_{-a}^{a} f(x) \, dx = 0\).
    \end{enumerate}
\end{thm}

\begin{defn} (séries de Fourier para funções pares/ímpares).
    \begin{enumerate}
        \item[(i)] A série de Fourier de uma função par \(f\) definida no intervalo \((-p,p)\) é a \textbf{série em cossenos}
        \begin{equation}
            f(x) = \frac{a_0}{2} + \sum_{n=1}^{\infty} a_n \cos\left(\frac{n\pi x}{p}\right),
        \end{equation}
        \noindent onde 
        \begin{align}
            a_0 &= \frac{2}{p} \int_{0}^{p} f(x) \, dx, \\
            a_n &= \frac{2}{p} \int_{0}^{p} f(x) \cos\left(\frac{n\pi x}{p}\right) \, dx.
        \end{align}
        \item[(ii)] A série de Fourier de uma função ímpar \(f\) definida no intervalo \((-p,p)\) é a \textbf{série em senos}
        \begin{equation}
            f(x) = \sum_{n=1}^{\infty} b_n \sin\left(\frac{n\pi x}{p}\right),
        \end{equation}
        \noindent onde
        \begin{equation}
            b_n = \frac{2}{p} \int_{0}^{p} f(x) \sin\left(\frac{n\pi x}{p}\right) \, dx.
        \end{equation}
    \end{enumerate}
\end{defn}

\subsection{Problema de Sturm-Liouville}

\renewcommand{\arraystretch}{1.15}
\begin{rem} (revisão de EDOs) \newline
    \begin{tabular}[c]{ll}
        \textbf{Equação de coeficientes constantes} & \textbf{Soluções gerais} \\
        \hline
        \(y'+\alpha y=0\) & \(y=ce^{-\alpha x}\) \\
        \(y''+\alpha^2 y=0,\quad \alpha>0\) & \(y=c_1\cos(\alpha x)+c_2\sin(\alpha x)\) \\
        \(y''-\alpha^2 y=0,\quad \alpha>0\) & \(\begin{cases} y=c_1e^{-\alpha x}+c_2e^{\alpha x} \\ y=c_1\cosh(\alpha x)+c_2\sinh(\alpha x) \end{cases}\) \\
        \\ \textbf{Equação de Cauchy-Euler} & \textbf{Soluções gerais, \(x>0\)} \\
        \hline
        \(x^2y''+ xy'- \alpha^2 y=0,\quad \alpha\geq 0\) & \(\begin{cases} y=c_1x^{-\alpha}+c_2x^{\alpha},\quad\alpha>0 \\ y=c_1+c_2\ln(x),\quad \alpha=0 \end{cases}\) \\
        \\ \textbf{Equação paramétrica de Bessel \(\nu=0\)} & \textbf{Solução geral, \(x>0\)} \\
        \hline
        \(x^2y''+ xy' + \alpha^2 x^2 y=0\) & \(y=c_1J_0(\alpha x)+c_2Y_0(\alpha x)\) \\
        \\ \textbf{Equação de Legendre \(n=0,1,2,\ldots\)} & \textbf{Soluções particulares são polinômios} \\
        \hline
        \((1-x^2)y'' - 2xy' + n(n+1)y=0\) & \(\begin{cases}
            y=P_0(x)=1,\\
            y=P_1(x)=x,\\
            y=P_2(x)=\frac{1}{2}(3x^2-1),\ldots
        \end{cases}\)
    \end{tabular}
\end{rem}

\begin{defn}
    Sejam \(p\), \(q\), \(r\) e \(r'\) funções reais contínuas em um intervalo \([a,b]\), e seja \(r(x)>0\) e \(p(x)>0\) para todo \(x \in [a,b]\). O problema de valor de contorno
    \begin{equation}
        \frac{d}{dx}\left[r(x)\frac{dy}{dx}\right] + q(x)y + \lambda p(x)y = 0, \quad a < x < b,
    \end{equation}
    \noindent sujeito às condições de contorno lineares e homogêneas
    \begin{align}
        A_1 y(a) + A_2 y'(a) &= 0, \\
        B_1 y(b) + B_2 y'(b) &= 0,
    \end{align}
    \noindent onde \(A_1\), \(A_2\), \(B_1\) e \(B_2\) são constantes tais que \(A_1^2 + A_2^2 \neq 0\) e \(B_1^2 + B_2^2 \neq 0\), é chamado de \emph{problema de Sturm-Liouville regular}. Os valores de \(\lambda\) para os quais existem soluções não-triviais são chamados de \emph{autovalores}, e as soluções correspondentes são chamadas de \emph{autovetores}.
\end{defn}

\begin{thm} (algumas propriedades do problema de Sturm-Liouville regular).
    \begin{enumerate}
        \item[\textup{(a)}] Existe um número infinito de autovalores reais que podem ser ordenados em uma sequência crescente \( \lambda_1 < \lambda_2 < \lambda_3 < \ldots \) tal que \( \lim_{n \to \infty} \lambda_n = +\infty \).
        \item[\textup{(b)}] A cada autovalor \( \lambda_n \) corresponde um único autovetor \( y_n(x) \), a menos de um fator multiplicativo constante.
        \item[\textup{(c)}] Autovetores correspondentes a autovalores distintos são linearmente independentes.
        \item[\textup{(d)}] O conjunto de autovetores correspondente ao conjunto de autovalores é ortogonal com respeito à função peso \( p(x) \) no intervalo \([a,b]\).
    \end{enumerate}
\end{thm}

\subsection{Séries de Fourier-Bessel e Séries de Fourier-Legendre}

\begin{defn}
    A série de \emph{Fourier-Bessel} de uma função \(f\) definida no intervalo \((0,b)\) é dada por
    \begin{enumerate}
        \item [(i)] \begin{gather}
            f(x) = \sum_{i=1}^{\infty} c_i J_n(\alpha_{i}x) \\
            c_i = \frac{2}{b^2 J_{n+1}^2(\alpha_i b)} \int_0^b x f(x) J_n(\alpha_i x) \, dx
        \end{gather}
        \noindent onde os \(\alpha_i\) são as raízes de \(J_n(\alpha b) = 0\).
        \item [(ii)] \begin{gather}
            f(x) = \sum_{i=1}^{\infty} c_i J_n(\alpha_{i}x) \\
            c_i = \frac{2\alpha_{i}^2}{(\alpha_{i}^2 b^2 - n^2 + h^2)J_n^2(\alpha_i b)} \int_0^b x f(x) J_n(\alpha_i x) \, dx
        \end{gather}
        \noindent onde os \(\alpha_i\) são as raízes de \(h J_n(\alpha b) + \alpha b J_n'(\alpha b) = 0\).
        \item [(iii)] \begin{gather}
            f(x) = c_1 + \sum_{i=1}^{\infty} c_i J_0(\alpha_{i}x) \\
            c_1 = \frac{2}{b^2} \int_0^b x f(x) \, dx, \quad c_i = \frac{2}{b^2 J_0^2(\alpha_i b)} \int_0^b x f(x) J_0(\alpha_i x) \, dx
        \end{gather}
        \noindent onde os \(\alpha_i\) são as raízes de \(J_0'(\alpha b) = 0\).
    \end{enumerate}
\end{defn}

\begin{prop} (convergência da série de Fourier-Bessel).
    Sejam \(f\) e \(f'\) contínuas por partes em \([0,b]\). Então para todo \(x\) no intervalo \((0,b)\), a série de Fourier-Bessel de \(f\) converge para \(f(x)\) em um ponto de continuidade. Em um ponto de descontinuidade, a série de Fourier-Bessel converge para a média
    \begin{equation*}
        \frac{f(x+) + f(x-)}{2},
    \end{equation*}
    \noindent onde \(f(x+)\) e \(f(x-)\) são os limites laterais de \(f\) em \(x\), respectivamente.
\end{prop}

\begin{defn}
    A série de \emph{Fourier-Legendre} de uma função \(f\) definida no intervalo \((-1,1)\) é dada por
    \begin{equation}
        f(x) = \sum_{n=0}^{\infty} c_n P_n(x),
    \end{equation}
    \noindent onde os coeficientes \(c_n\) são dados por
    \begin{equation}
        c_n = \frac{2n+1}{2} \int_{-1}^{1} f(x) P_n(x) \, dx.
\end{equation}
\end{defn}

\begin{prop} (convergência da série de Fourier-Legendre).
    Sejam \(f\) e \(f'\) contínuas por partes em \([-1,1]\). Então para todo \(x\) no intervalo \((-1,1)\), a série de Fourier-Legendre de \(f\) converge para \(f(x)\) em um ponto de continuidade. Em um ponto de descontinuidade, a série de Fourier-Legendre converge para a média
    \begin{equation*}
        \frac{f(x+) + f(x-)}{2},
    \end{equation*}
    \noindent onde \(f(x+)\) e \(f(x-)\) são os limites laterais de \(f\) em \(x\), respectivamente.
\end{prop}

\section{Problemas de Valor de Contorno em Coordenadas Retangulares}

\subsection{Equações Diferenciais Parciais Separáveis}

\begin{thm} (princípio da superposição).
    Se \(u_1, u_2, \ldots, u_k\) são soluções de uma EDP linear homogênea, então a combinação linear
    \begin{equation}
        u = c_1 u_1 + c_2 u_2 + \cdots + c_k u_k,
    \end{equation}
    \noindent onde \(c_1, c_2, \ldots, c_k\) são constantes, também é uma solução da EDP.
\end{thm}

\begin{rem}
    Vamos assumir que quando tivermos um conjunto infinito \(u_1, u_2, \ldots\) de soluções de uma equação linear homogênea, podemos construir outra solução \(u\) como uma série infinita
    \begin{equation}
        u = \sum_{k=1}^{\infty} c_k u_k,
    \end{equation}
    \noindent onde \(c_1, c_2, \ldots\) são constantes.
\end{rem}

\begin{defn}
    A equação diferencial parcial linear de segunda ordem
    \begin{equation}
        A \frac{\partial^2 u}{\partial x^2} + B \frac{\partial^2 u}{\partial x \partial y} + C \frac{\partial^2 u}{\partial y^2} + D \frac{\partial u}{\partial x} + E \frac{\partial u}{\partial y} + F u = G,
    \end{equation}
    \noindent onde \(A\), \(B\), \(C\), \(D\), \(E\), \(F\) e \(G\) são constantes reais, é classificada como:
    \begin{itemize}
        \item \emph{hiperbólica} se \(B^2 - 4AC > 0\),
        \item \emph{parabólica} se \(B^2 - 4AC = 0\),
        \item \emph{elíptica} se \(B^2 - 4AC < 0\).
    \end{itemize}
\end{defn}

\begin{defn} (equações clássicas).
    \begin{align}
        &k u_{xx} = u_t, \quad k>0 \label{eq:calor-1d} \\
        &a^2 u_{xx} = u_{tt} \label{eq:onda-1d} \\
        &\nabla^2 u = 0 \label{eq:laplace-2d}
    \end{align}
    \noindent onde \(\nabla^2 = \partial^2/\partial x^2 + \partial^2/\partial y^2\).
    As equações \eqref{eq:calor-1d}, \eqref{eq:onda-1d} e \eqref{eq:laplace-2d} são chamadas, respectivamente, de \emph{equação do calor unidimensional}, \emph{equação da onda unidimensional} e \emph{equação de Laplace bidimensional}, onde a constante \(k\) é o coeficiente de difusão térmica e \(a\) é a velocidade da onda no meio.
    Note que a equação do calor é parabólica, a equação da onda é hiperbólica e a equação de Laplace é elíptica.
\end{defn}

\begin{rem}
    Em geral, existem três tipos de condições de contorno:
    \begin{itemize}
        \item[(i)] Condições de Dirichlet: especificam o valor da função na fronteira do domínio (ex.: \(u=0\) na fronteira).
        \item[(ii)] Condições de Neumann: especificam o valor da derivada normal da função na fronteira do domínio (ex.: \(\partial u/\partial n=0\) na fronteira).
        \item[(iii)] Condições de Robin: combinam condições de Dirichlet e Neumann em diferentes partes da fronteira do domínio (ex.: \(\partial u/\partial n + hu = g\) na fronteira, onde \(h\) é uma constante e \(g\) é uma função definida na fronteira).
    \end{itemize}
\end{rem}

\section{Problemas de Valor de Contorno em Outros Sistemas de Coordenadas}

\subsection{Coordenadas Polares}
\subsection{Coordenadas Polares e Cilíndricas}
\subsection{Coordenadas Esféricas}

\section{Transformadas Integrais}

\subsection{Transformada de Laplace}

\begin{defn}
    A \emph{transformada de Laplace} de uma função \(f(t)\), definida para \(t \geq 0\), é a função \(\mathcal{L}\{f(t)\} = F(s)\) dada por
    \begin{equation}
        F(s) = \int_{0}^{\infty} f(t) e^{-st} \, dt,
    \end{equation}
    \noindent onde \(s\) é um número complexo.
\end{defn}

\begin{prop} (transformada de Laplace de derivadas parciais).
    As transformadas das derivadas parciais \(\partial u/\partial t\) e \(\partial^2 u/\partial t^2\) são dadas por
    \begin{align}
        \mathcal{L}\left\{\frac{\partial u}{\partial t}\right\} &= s U(x,s) - u(x,0), \\
        \mathcal{L}\left\{\frac{\partial^2 u}{\partial t^2}\right\} &= s^2 U(x,s) - s u(x,0) - \frac{\partial u}{\partial t}(x,0),
    \end{align}
    \noindent onde \(U(x,s) = \mathcal{L}\{u(x,t)\}\).

    Como estamos transformando em relação a \(t\), assumimos que é possível intercambiar integração e diferenciação na transformada de \(\partial^2u/\partial x^2\), ou seja,
    \begin{equation}
        \mathcal{L}\left\{\frac{\partial^2 u}{\partial x^2}\right\} = \frac{\partial^2}{\partial x^2} \mathcal{L}\{u(x,t)\} = \frac{d^2 U}{d x^2}(x,s).
    \end{equation}
\end{prop}

\subsection{Integral de Fourier}

\begin{defn}
    A \emph{integral de Fourier} de uma função \(f\) definida no intervalo \(\left(-\infty, \infty\right)\) é dada por
    \begin{equation}\label{eq:integral-fourier}
        f(x) = \frac{1}{\pi} \int_{0}^{\infty} \left[A(\alpha) \cos(\alpha x) + B(\alpha) \sin(\alpha x)\right] d\alpha,
    \end{equation}
    \noindent onde
    \begin{align}
        A(\alpha) &= \int_{-\infty}^{\infty} f(x) \cos(\alpha x) \, dx, \\
        B(\alpha) &= \int_{-\infty}^{\infty} f(x) \sin(\alpha x) \, dx.
    \end{align}
\end{defn}

\begin{thm} (condições para convergência da integral de Fourier).
    Seja \(f\) e \(f'\) contínuas por partes em todo intervalo finito e seja \(f\) absolutamente integrável em \((-\infty, \infty)\). Então, a integral de Fourier de \(f\) converge para \(f(x)\) em um ponto de continuidade. Em um ponto de descontinuidade, a integral de Fourier converge para a média
    \begin{equation}
        \frac{f(x+) + f(x-)}{2},
    \end{equation}
    \noindent onde \(f(x+)\) e \(f(x-)\) são os limites laterais de \(f\) em \(x\), respectivamente.
\end{thm}

\begin{defn} (integrais de Fourier em senos e cossenos).
    \begin{enumerate}
        \item[(i)] A integral de Fourier de uma função par \(f\) definida no intervalo \(\left(-\infty, \infty\right)\) é a \textbf{integral em cossenos}
        \begin{equation}
            f(x) = \frac{2}{\pi} \int_{0}^{\infty} A(\alpha) \cos(\alpha x) \, d\alpha,
        \end{equation}
        \noindent onde
        \begin{align}
            A(\alpha) &= \int_{0}^{\infty} f(x) \cos(\alpha x) \, dx.
        \end{align}
        \item[(ii)] A integral de Fourier de uma função ímpar \(f\) definida no intervalo \(\left(-\infty, \infty\right)\) é a \textbf{integral em senos}
        \begin{equation}
            f(x) = \frac{2}{\pi} \int_{0}^{\infty} B(\alpha) \sin(\alpha x) \, d\alpha,
        \end{equation}
        \noindent onde
        \begin{align}
            B(\alpha) &= \int_{0}^{\infty} f(x) \sin(\alpha x) \, dx.
        \end{align}
    \end{enumerate}
\end{defn}

\begin{rem}
    A integral de Fourier \eqref{eq:integral-fourier} também possui a forma complexa, ou forma exponencial, que é análoga à forma complexa da série de Fourier \eqref{eq:serie-fourier-complexa}.
    \begin{equation}\label{eq:integral-fourier-complexa}
        f(x) = \frac{1}{2\pi} \int_{-\infty}^{\infty} C(\alpha) e^{-i\alpha x} \, d\alpha,
    \end{equation}
    \noindent onde
    \begin{equation}
        C(\alpha) = \int_{-\infty}^{\infty} f(x) e^{i\alpha x} \, dx.
    \end{equation}
\end{rem}

\subsection{Transformada de Fourier}

\begin{defn} (pares de transformadas de Fourier).
    \begin{enumerate}
        \item[(i)] Transformada de Fourier: \begin{equation}
            \widehat{f}(\alpha) = \mathcal{F}\{f(x)\} = \int_{-\infty}^{\infty} f(x) e^{i\alpha x} \, dx.
        \end{equation}
        \item[] Transformada inversa de Fourier: \begin{equation}
            f(x) = \mathcal{F}^{-1}\{\widehat{f}(\alpha)\} = \frac{1}{2\pi} \int_{-\infty}^{\infty} \widehat{f}(\alpha) e^{-i\alpha x} \, d\alpha.
        \end{equation}
        \item[(ii)] Transformada de Fourier em senos: \begin{equation}
            \widehat{f}_s(\alpha) = \mathcal{F}_s\{f(x)\} = \int_{0}^{\infty} f(x) \sin(\alpha x) \, dx.
        \end{equation}
        \item[] Transformada inversa de Fourier em senos: \begin{equation}
            f(x) = \mathcal{F}_s^{-1}\{\widehat{f}_s(\alpha)\} = \frac{2}{\pi} \int_{0}^{\infty} \widehat{f}_s(\alpha) \sin(\alpha x) \, d\alpha.
        \end{equation}
        \item[(iii)] Transformada de Fourier em cossenos: \begin{equation}
            \widehat{f}_c(\alpha) = \mathcal{F}_c\{f(x)\} = \int_{0}^{\infty} f(x) \cos(\alpha x) \, dx.
        \end{equation}
        \item[] Transformada inversa de Fourier em cossenos: \begin{equation}
            f(x) = \mathcal{F}_c^{-1}\{\widehat{f}_c(\alpha)\} = \frac{2}{\pi} \int_{0}^{\infty} \widehat{f}_c(\alpha) \cos(\alpha x) \, d\alpha.
        \end{equation}
    \end{enumerate}
\end{defn}

\begin{rem}
    A \emph{transformada de Fourier} \(\mathcal{F}\) é um operador linear que transforma uma função \(f(x)\) em uma função \(\widehat{f}(\alpha)\) no domínio da frequência. A \emph{transformada inversa de Fourier} \(\mathcal{F}^{-1}\) recupera a função original \(f(x)\) a partir de sua transformada \(\widehat{f}(\alpha)\).
    Em outras palavras, \(\mathcal{F}\{f(x)\}\) é a projeção de \(f(x)\) sobre as funções base \(e^{i\alpha x}\) do espaço de funções, enquanto \(\mathcal{F}^{-1}\{\widehat{f}(\alpha)\}\) reconstrói \(f(x)\) a partir dessas projeções.
\end{rem}

\begin{res} (algumas propriedades das transformadas de Fourier).
    Suponha que \(f\) é uma função contínua e absolutamente integrável no intervalo \((-\infty, \infty)\) e que \(f'\) seja contínua por partes em todo intervalo finito. Se \(f(x) \to 0\) quando \(|x| \to \infty\), então:
    \begin{equation}
        \mathcal{F}\{f'(x)\} = i\alpha \widehat{f}(\alpha).
    \end{equation}

    Se \(f\) e \(f'\) são contínuas, \(f\) é absolutamente integrável em \((0, \infty)\), e \(f''\) é contínua por partes em todo intervalo finito, então:
    \begin{enumerate}
        \item [\textup{(a)}] \(\mathcal{F}_s\{f''(x)\} = -\alpha^2 \widehat{f}_c(\alpha) + \alpha f(0)\).
        \item [\textup{(b)}] \(\mathcal{F}_c\{f''(x)\} = \alpha^2 \widehat{f}_s(\alpha) - f'(0)\).
    \end{enumerate}
\end{res}

\end{document}